% =========================================================================
% 二、问题分析
% =========================================================================

\section{问题分析}

四问共享同一物理系统与同一储能约束，差别在于\textbf{信息条件与结算规则}。因此本文的问题分析统一采用``新约束或信息变化 $\to$ 为什么原方法不足 $\to$ 本文作何修改''的三步框架，把定量结果与对照留到各问的结果分析，不在前置分析中提前宣布结论。技术路线随之从单日线性规划升级为``预测—计划—执行''三层滚动框架（\S\ref{sec:q2-frame}），再扩展为``计划—调整—执行''三层结构（\S\ref{sec:q3-frame}），最后在价格侧引入不确定性并延展优化时域（\S\ref{sec:q4-model}）。

\subsection{问题一：时段转移与富余吸收}

\textbf{信息与约束变化。}当天负载与光伏均视为已知，决策退化为确定性优化；但储能把各区间耦合起来——当前区间的充放电通过储电量影响后续所有区间。

\textbf{原方法的不足。}逐区间贪心无法处理上述跨时段耦合；而物理上储能不可同时充放电（$c_td_t=0$），直接引入该约束将使模型非线性，失去求全局最优的条件。

\textbf{本文的修改。}建立全天联合的线性规划（\S\ref{sec:q1-constraint}）；并在 \S\ref{sec:q1-relax} 证明互斥约束的线性松弛在最优值意义下精确，模型保持为可求全局最优的线性规划。

\begin{figure}[H]
  \centering
  \includegraphics[width=0.82\textwidth]{p1_fig1.pdf}
  \caption{代表日输入特征：小区负载、光伏预测功率与电价曲线}
  \label{fig:q1-input}
\end{figure}

\subsection{问题二：不对称代价下的序贯决策}

\textbf{信息与约束变化。}计划只能基于预测，实测偏差必须按 5 倍电价紧急补足。这一惩罚使代价结构高度不对称：少计划 1 kWh 的期望代价远高于多计划 1 kWh 的代价。

\textbf{原方法的不足。}逐日实测净负载同时含有星期周期与随季节缓慢漂移的近期水平两层可预测结构。只用近期水平会平滑掉星期形状，只用代表日曲线又跟不上季节漂移；此外，任何用``当天实测''改善``当天计划''的做法都不可实施，需要用一条明确的信息纪律把这种泄漏排除掉。

\textbf{本文的修改。}预测层显式分解为``近期水平＋星期偏差''两层，再叠加用于吸收残余随机性的分位裕量（\S\ref{sec:q2-forecast}）；计划层保持线性规划（\S\ref{sec:q2-lp}）；执行层采用完全因果的尽限充放电规则，在净缺口出现时即时释放可用电量（\S\ref{sec:q2-exec}）；参数按逐月 walk-forward 协议选取（\S\ref{sec:q2-walkforward}）。

\subsection{问题三：信息增加与调整机会的价值辨析}
\label{sec:q3-analysis}

\textbf{信息与约束变化。}题面并未要求``利用更多预报提高预测精度''，而是要求分析\textbf{是否需要}引入其他时刻的预报。除新增 6:00、12:00、18:00 三个预报时刻外，还允许对这些时刻之后的购电量\textbf{重新约定}，代价是相对 0:00 计划的偏离要按双向违约价付费。

\textbf{原方法的不足。}问题二的一次性决策无法利用日内信息，这是明确的改进方向；但``信息更多''与``费用更低''之间不存在必然的对应关系，需要分别计量两条作用渠道：其一，调整机会本身有价值——白天若发现净负载低于预期，可下调合同并按违约价结算；其二，调整本身有摩擦——双向违约价使``偏离 0:00 计划''成为独立成本项，且该成本与预测精度直接耦合，预测越激进地偏离基准，重定合同的动作越大。此外，两项改进可能部分替代：已有第二次决策机会时，0:00 预报的边际价值会被削弱。

\textbf{本文的修改。}插入\textbf{调整层}，形成``计划—调整—执行''三层结构（\S\ref{sec:q3-frame}、\S\ref{sec:q3-adjust}）；并采用 $2\times2$ 对照（\S\ref{sec:q3-matrix}）而非简单相减来计量``信息''与``调整机会''各自的收益及其重叠。

\subsection{问题四：价格侧不确定性的来源辨析与信息边界}
\label{sec:q4-analysis}

\textbf{信息与约束变化。}问题二、问题三的不确定性全部位于负载与光伏一侧，问题四新增了价格一侧：电价由已知量变为逐区间的随机量，且按本文解释随区间推进逐步揭示。

\textbf{原方法的不足。}两处需要修改：其一，问题二的均一分位裕量对所有时段取同一分位水平，未反映``缺口发生在高价时段时单位代价更高''这一事实；其二，日循环边界限制的是\textbf{计划中的单日净库存变化}，在固定电价下这一限制尚不显著，而在实时价格下，低价夜间若不能把电量安排在更晚的时段使用，低价优势无法充分利用。同时，两张不确定性谁更重要不能凭直觉判断，需要先分解价格的统计结构，再分别标定两侧信息的经济价值。

\textbf{本文的修改。}先作价格结构分解与价格缩放不变性分析（\S\ref{sec:q4-struct}、\S\ref{sec:q4-scaling}），据此把价格预测的目标由``全曲线''压缩到``归一化形状''；模型层面引入按价格形状加权的风险裕量（式~\eqref{eq:q4-prm}），并把优化时域延展到次日 24:00（\S\ref{sec:q4-horizon}），同时在信息侧严格保持因果（次日只用持久性预测）；结果以同政策对照、终端条件一致性与放松下界参照分别检验（\S\ref{sec:q4-ablation}、\S\ref{sec:q4-regret}）。
